% ============================================================
% Preamble — Machine Learning Notes
% ============================================================

% --- Fonts (XeTeX/tectonic, unicode math) ---------------------------------
\usepackage{amsmath}
\usepackage{fontspec}
\usepackage{unicode-math}
% (STIX Two Text on macOS is a variable font: XeTeX cannot find its bold face)
\setmainfont{Charter}
\setsansfont{Helvetica Neue}
\setmonofont{Menlo}[Scale=0.85]
\setmathfont{STIX Two Math}
\AtBeginDocument{\providecommand{\blacksquare}{\mdlgblksquare}}
% glyphs missing from Charter: taken from the math font
\usepackage{newunicodechar}
\newunicodechar{→}{\ensuremath{\rightarrow}}
\newunicodechar{−}{\ensuremath{-}}
\newunicodechar{×}{\ensuremath{\times}}
\newunicodechar{₁}{\ensuremath{_1}}
\newunicodechar{²}{\ensuremath{^2}}
\newunicodechar{±}{\ensuremath{\pm}}

% --- English language (chapter/contents names, \today, quoting) -----------
\usepackage{polyglossia}
\setmainlanguage{english}

% --- Page geometry ----------------------------------------------------------
\usepackage[a4paper,margin=2.6cm,headheight=14pt]{geometry}
% book class defaults to \flushbottom, which stretches vertical glue between
% paragraphs to fill every page -- with many figures/boxes this produces huge
% random gaps. Let pages run slightly short instead.
\raggedbottom
\setlength{\parskip}{3pt plus 1pt}
% a bit of extra stretch avoids overfull lines with long unhyphenated terms
\setlength{\emergencystretch}{3em}

% --- Colors -----------------------------------------------------------------
\usepackage{xcolor}
\definecolor{inkblue}{HTML}{1B3A6B}
\definecolor{accentblue}{HTML}{2F5FA8}
\definecolor{defback}{HTML}{EDF2FB}
\definecolor{defframe}{HTML}{2F5FA8}
\definecolor{thmback}{HTML}{F1EDFA}
\definecolor{thmframe}{HTML}{6A4C93}
\definecolor{exback}{HTML}{EAF7EE}
\definecolor{exframe}{HTML}{2E8B57}
\definecolor{noteback}{HTML}{F2F2F2}
\definecolor{noteframe}{HTML}{6B6B6B}
\definecolor{tipback}{HTML}{E7F6F4}
\definecolor{tipframe}{HTML}{1E8C7C}
\definecolor{warnback}{HTML}{FDF1E4}
\definecolor{warnframe}{HTML}{C0722A}
\definecolor{absback}{HTML}{EFEFF5}
\definecolor{absframe}{HTML}{4A4A68}
\definecolor{qback}{HTML}{FBF6E3}
\definecolor{qframe}{HTML}{A8861B}

% --- Boxes for Obsidian callouts --------------------------------------------
% \newcallout{env}{label}{back}{frame}: the box title is "label — title",
% or just "label" when the callout has no title.
\usepackage{etoolbox}
\usepackage[most]{tcolorbox}
\tcbuselibrary{skins,breakable}
\newcommand{\newcallout}[4]{%
  \newtcolorbox{#1}[1]{%
    enhanced, breakable, colback=#3, colframe=#4,
    boxrule=0.5pt, left=6pt, right=6pt, top=4pt, bottom=4pt,
    arc=2pt, before skip=8pt, after skip=8pt,
    fonttitle=\bfseries, coltitle=white, colbacktitle=#4,
    title={#2\ifblank{##1}{}{ — ##1}}}}
\newcallout{boxdefinition}{Definition}{defback}{defframe}
\newcallout{boxtheorem}{Theorem}{thmback}{thmframe}
\newcallout{boxexample}{Example}{exback}{exframe}
\newcallout{boxnote}{Note}{noteback}{noteframe}
\newcallout{boxtip}{Tip}{tipback}{tipframe}
\newcallout{boxwarning}{Warning}{warnback}{warnframe}
\newcallout{boxabstract}{Summary}{absback}{absframe}
\newcallout{boxquestion}{Questions}{qback}{qframe}

% --- Code blocks --------------------------------------------------------------
\usepackage{fvextra}
\fvset{breaklines=true,breakanywhere=true,fontsize=\small,frame=leftline,
       rulecolor=\color{accentblue!50},framesep=6pt,xleftmargin=8pt}

% --- Tables / figures ---------------------------------------------------------
\usepackage{longtable}
\usepackage{booktabs}
\usepackage{array}
\usepackage{calc}
\usepackage{graphicx}
\usepackage{float}
\usepackage{caption}
% floats: fill pages near the insertion point, avoiding half-empty pages
\renewcommand{\topfraction}{0.85}
\renewcommand{\bottomfraction}{0.7}
\renewcommand{\textfraction}{0.1}
\renewcommand{\floatpagefraction}{0.75}
\setcounter{topnumber}{3}
\setcounter{bottomnumber}{2}
\setcounter{totalnumber}{5}
\setlength{\textfloatsep}{10pt plus 2pt minus 4pt}
\setlength{\floatsep}{8pt plus 2pt minus 2pt}
\setlength{\intextsep}{8pt plus 2pt minus 2pt}
\captionsetup{font=small,labelfont=bf,margin=1cm,skip=5pt}
% pandoc helpers
\providecommand{\tightlist}{\setlength{\itemsep}{0pt}\setlength{\parskip}{0pt}}
\newcommand{\real}[1]{#1}
\newcounter{none} % used by pandoc for longtables without caption

% --- Lists ----------------------------------------------------------------
\usepackage{enumitem}
\setlist{itemsep=2pt, topsep=4pt}

% --- Chapter / section styling ------------------------------------------------
\usepackage{titlesec}
\titleformat{\chapter}[display]
  {\normalfont\huge\bfseries\color{inkblue}}
  {\Large\color{accentblue}\chaptertitlename\ \thechapter}
  {8pt}{\Huge}
\titlespacing*{\chapter}{0pt}{0pt}{28pt}

\titleformat{\section}
  {\normalfont\Large\bfseries\color{inkblue}}{\thesection}{0.6em}{}
\titleformat{\subsection}
  {\normalfont\large\bfseries\color{accentblue}}{\thesubsection}{0.6em}{}
\titleformat{\subsubsection}
  {\normalfont\normalsize\bfseries\itshape\color{accentblue}}{\thesubsubsection}{0.6em}{}
\setcounter{secnumdepth}{2}
\setcounter{tocdepth}{2}

% --- Header / footer ----------------------------------------------------------
\usepackage{fancyhdr}
\pagestyle{fancy}
\fancyhf{}
\fancyhead[R]{\thepage}
\fancyhead[L]{\nouppercase{\leftmark}}
\renewcommand{\headrulewidth}{0.4pt}

% --- Chapter metadata line under chapter title --------------------------------
\newcommand{\lecturemeta}[2]{%
  \begingroup
  \color{gray!70!black}\small
  #1\\
  \textit{Source: #2}
  \endgroup
  \bigskip\par
}

% --- Hyperref (load last) ----------------------------------------------------
\usepackage[bookmarksnumbered,hidelinks]{hyperref}
\hypersetup{
  colorlinks=true,
  linkcolor=accentblue,
  citecolor=accentblue,
  urlcolor=accentblue,
  pdftitle={Machine Learning — Notes},
  pdfauthor={Fabio Piscitelli}
}
\usepackage{xurl}
